\documentclass[11pt]{article}
\usepackage{amsmath,amssymb,amsthm,mathtools}
\usepackage[margin=1in]{geometry}
\newtheorem{theorem}{Theorem}
\newtheorem{proposition}[theorem]{Proposition}
\newtheorem{lemma}[theorem]{Lemma}
\newtheorem{corollary}[theorem]{Corollary}
\theoremstyle{definition}
\newtheorem{remark}[theorem]{Remark}
\theoremstyle{remark}
\newtheorem*{verdict}{Verdict}
\DeclareMathOperator{\Res}{Res}
\newcommand{\PsiC}{\Psi^{\mathrm{C}}}
\title{Q553: the one-head residue ties are not pole-groupings.\\
A test-function basis, a reflection duality, and an alternant\\
expansion that collapses every head to length one}
\author{Clio}
\date{2026-10-22 (prove, cycle 2)}

\begin{document}
\maketitle

\begin{abstract}
Q553 asked whether the $u\to\infty$ tie family is exactly the family of admissible
pole-groupings of the single rational function $R$ underlying my one-head residue
machine. The answer is \textbf{no}, and for a structural reason: the residue theorem on
$\mathbb{P}^1$ is \emph{one} linear relation, so all $2^{N}-2$ groupings of a fixed $R$
are rebracketings of a single identity (Theorem~\ref{thm:one}). The correct index is not
a grouping but a \emph{test function}, and $\{z^a\}_{a\in\mathbb{Z}}$ is a basis of the
resulting family (\S\ref{sec:basis}), each member with a three-case closed form
(Theorem~\ref{thm:table}). The family carries a reflection duality
$I_a(1/w)=-I_{-2-a}(w)$ whose unique fixed point $a=-1$ is precisely the $\lambda=(2,3)$
tie (Theorem~\ref{thm:dual}). Separately, and by an argument needing no residues at all,
every head length collapses to length one (Theorem~\ref{thm:collapse}); grouping the
head sum by its underlying subset exhibits $\PsiC_{(1^r),\theta}$ as a weighted sum of
\emph{alternants} (Theorem~\ref{thm:alt}), which proves a $\rho$-shifted vanishing law:
$\PsiC_{(1^r),\theta}=0$ whenever $\theta_j-\theta_k=k-j$ for some $j<k$
(Theorem~\ref{thm:vanish}). Every claim is verified in exact rational arithmetic; the
counts are in \S\ref{sec:verif}. The session's declared first task (T0) died on
provenance, not mathematics; \S\ref{sec:T0} records that honestly.
\end{abstract}

\section{Setup}

Fix a finite index set $S\subset\mathbb{Z}_{\ge1}$, write $n=|S|$, and let
$w=(w_k)_{k\ge1}$ be indeterminates specialised to nonzero scalars. Put
\[
  e_S=\prod_{m\in S}w_m,\qquad
  P(z)=\prod_{m\in S}(w_mz-1),\qquad
  D(z)=\prod_{m\in S}(z-w_m),\qquad
  Q_s(z)=\prod_{k=1}^{s-1}(z-w_k),
\]
and for $s\ge1$, $t\in\mathbb{Z}$,
\begin{equation}\label{eq:R}
  R_{s,t}(z)\;=\;z^{\,t-s}\,Q_s(z)\,\frac{P(z)}{D(z)}.
\end{equation}
Let $\Phi$ be the Lagrange functional
\begin{equation}\label{eq:Phi}
  \Phi[g]\;=\;\sum_{m\in S}\frac{g(w_m)}{\prod_{t\in S,\,t\neq m}(w_m-w_t)} .
\end{equation}
The one-head specialisation of my $\PsiC$ (the $\Psi$ of \emph{my} registry --- not
Ikeda--Naruse's GKM $\Psi_n$, not Rupel's $\Psi_i$; see Remark~\ref{rem:psi}) satisfies,
at trust \texttt{proved} ($75/75$ exact, \texttt{reviews/scripts-2026-10-22/residue\_lemma.py},
node \texttt{clio-one-head-Psi-residue-closed-form}),
\begin{equation}\label{eq:proved}
  e_S\cdot\PsiC_{(s),t}(S)\;=\;-\Res_{z=0}R_{s,t}\;-\;\Res_{z=\infty}R_{s,t}.
\end{equation}

\begin{remark}\label{rem:psi}
Three different objects named $\Psi$ are live in the surrounding conversation:
Ikeda--Naruse's GKM ring $\Psi_n$, Rick's $\Psi_{\kappa,\theta}$, and Rupel's $\Psi_i$.
Throughout this note $\PsiC$ denotes \emph{mine}, defined by the head-permutation sum
\eqref{eq:Psidef} below. Every $\Psi$ in this document carries the superscript.
\end{remark}

Explicitly, for $\kappa\in\mathbb{Z}_{\ge1}^r$ and $\theta\in\mathbb{Z}^r$, with
$\rho(\kappa,X,w)=\prod_{k<\kappa}(1-w_k/X)$,
\begin{equation}\label{eq:Psidef}
  \PsiC_{\kappa,\theta}(M)
  =\sum_{\substack{h=(h_1,\dots,h_r)\\ \text{distinct in }M}}
   \;\prod_{j=1}^{r}\Bigl(1-X_j^{-2}\Bigr)\rho(\kappa_j,X_j,w)X_j^{\theta_j}
   \prod_{j<k}\frac{1-\frac{1}{X_jX_k}}{1-\frac{X_j}{X_k}}
   \prod_{\substack{t\in M\\ t\notin h}}\prod_{j=1}^{r}\frac{1-\frac{1}{w_tX_j}}{1-\frac{w_t}{X_j}},
\end{equation}
where $X_j=w_{h_j}$.

\section{Q553 as asked: the groupings carry one identity, not a family}

\begin{theorem}\label{thm:one}
Let $R$ be any rational function on $\mathbb{P}^1$ and let
$\mathcal{N}=\{p\in\mathbb{P}^1:\Res_pR\neq0\}$, $N=|\mathcal{N}|$. For each proper
nonempty $A\subsetneq\mathcal{N}$ the ``grouping identity''
$\sum_{p\in A}\Res_pR=-\sum_{p\in\mathcal{N}\setminus A}\Res_pR$ holds, and all
$2^{N}-2$ of them are equivalent to the single relation
$\sum_{p\in\mathcal{N}}\Res_pR=0$. Hence the space of relations they span on the residue
vector has rank $1$.
\end{theorem}

\begin{proof}
The residue theorem for the rational $1$-form $R\,dz$ on $\mathbb{P}^1$ gives
$\sum_{p\in\mathbb{P}^1}\Res_pR=0$, and residues vanish off $\mathcal{N}$, so
$\sum_{p\in\mathcal{N}}\Res_pR=0$. For any $A$, adding
$\sum_{p\in\mathcal{N}\setminus A}\Res_pR$ to both sides of the grouping identity
returns exactly this relation, and subtracting it from the relation returns the grouping
identity. The two are therefore equivalent, for every $A$. A single linear relation spans
a rank-$1$ space.
\end{proof}

\begin{verdict}
\textbf{Q553: NO.} A tie family of cardinality greater than one cannot be indexed by the
groupings of a single $R$, because those groupings carry one identity between them. The
brief's instruction to ``enumerate the $2^3-2=6$'' identities at $|S|=2$ was counting
bracketings of one relation, not identities. Note also that the premise contains a second,
independent error: the index set is not the \emph{pole} set (Theorem~\ref{thm:carrier}).
\end{verdict}

\begin{theorem}[the residue-carrying set is not the pole set]\label{thm:carrier}
Assume the $w_k$, for $k\in S\cup\{1,\dots,s-1\}$, are distinct, nonzero, and satisfy
$w_jw_m\neq1$ for all $j,m\in S$. Then
\[
  \mathcal{N}(S,s,t)\;\subseteq\;
  \{\,w_m: m\in S,\ m\ge s\,\}\;\cup\;\{\,0 \text{ if } t<s\,\}\;\cup\;\{\,\infty \text{ if } t\ge0\,\},
\]
with equality for generic $w$; the memberships of $w_m$ $(m\ge s)$, of $0$ at $t=s-1$,
and of $\infty$ at $t=0$ hold unconditionally.
\end{theorem}

\begin{proof}
\emph{At $z=0$.} $Q_s(0)=\prod_{k<s}(-w_k)\neq0$, $P(0)=(-1)^n\neq0$ and
$D(0)=(-1)^ne_S\neq0$, so $Q_sP/D$ is holomorphic and nonzero at $0$ and
$\operatorname{ord}_0R_{s,t}=t-s$ exactly. A nonzero residue at $0$ requires a $z^{-1}$
term, hence $t-s\le-1$, i.e.\ $t<s$. For $t=s-1$ the residue equals $(Q_sP/D)(0)\neq0$;
for $t<s-1$ it is a higher Taylor coefficient of $Q_sP/D$ at $0$, nonzero off a proper
subvariety.

\emph{At $z=\infty$.} The leading behaviour is
$z^{t-s}\cdot z^{s-1}\cdot e_Sz^{n}/z^{n}=e_Sz^{t-1}$, so
$R_{s,t}(z)=e_Sz^{t-1}\sum_{k\ge0}\gamma_kz^{-k}$ with $\gamma_0=1$, and
$\Res_\infty R_{s,t}=-e_S\gamma_t$. This requires $t\ge0$; at $t=0$ it equals
$-e_S\neq0$ unconditionally, and for $t\ge1$ it is $-e_S\gamma_t$, nonzero generically.
Observe that at $t=0$ the point $\infty$ is a \emph{zero} of $R_{s,t}$, not a pole,
yet carries the residue $-e_S$: this is the off-by-one in the premise.

\emph{At $z=w_m$, $m\in S$.} If $m\le s-1$ then $Q_s$ contains the factor $(z-w_m)$,
which cancels the simple zero of $D$ there, so $R_{s,t}$ is regular at $w_m$ and the
residue is $0$. If $m\ge s$ then $D$ has a simple zero at $w_m$ and $Q_s(w_m)\neq0$ by
distinctness, so the pole is simple with residue
$w_m^{t-s}Q_s(w_m)P(w_m)/\prod_{t'\neq m}(w_m-w_{t'})$, and
$P(w_m)=\prod_{j\in S}(w_jw_m-1)\neq0$ by hypothesis.
\end{proof}

\begin{remark}
Theorem~\ref{thm:carrier} is consistent with \eqref{eq:Psidef}: the head factor
$\rho(s,X,w)=\prod_{k<s}(1-w_k/X)$ vanishes at $X=w_k$ for $k<s$, so $\PsiC_{(s),t}$
drops exactly the heads $h<s$ that $Q_s$ cancels. The cancellation is visible on both
sides of \eqref{eq:proved}.
\end{remark}

\section{The right index: test functions}\label{sec:basis}

\begin{lemma}\label{lem:Phi}
For every $g\in\mathbb{C}[z,z^{-1}]$,
$\displaystyle \Phi[g]=-\Res_{z=0}\bigl[g/D\bigr]-\Res_{z=\infty}\bigl[g/D\bigr]$.
\end{lemma}

\begin{proof}
$D$ has simple zeros at the distinct $w_m\neq0$, at which $g$ is regular, so
$\Res_{w_m}[g/D]=g(w_m)/D'(w_m)=g(w_m)/\prod_{t\neq m}(w_m-w_t)$; summing gives
$\Phi[g]=\sum_{m\in S}\Res_{w_m}[g/D]$. The only other possible poles of the rational
function $g/D$ are $0$ and $\infty$, so the residue theorem on $\mathbb{P}^1$ gives the
claim.
\end{proof}

Write $I_a:=\Phi[z^aP]$ for $a\in\mathbb{Z}$. Since $f\mapsto\Phi[fP]$ is linear and
$R_{s,t}=fP/D$ with $f=z^{t-s}Q_s$, Lemma~\ref{lem:Phi} says the identity family is
indexed by the \emph{test function} $f$, one relation per $f$; and $\{z^a\}_{a\in\mathbb{Z}}$
is a basis of $\mathbb{C}[z,z^{-1}]$. Taking $s=1$ (so $Q_1=1$) already realises every
monomial, whence
\begin{equation}\label{eq:Ia}
  I_a=e_S\cdot\PsiC_{(1),a+1}(S),
\end{equation}
and the $(s,t)$-family is a triangular re-basing of the monomial family, with unit
diagonal (the coefficient of $z^{t-1}$ in $z^{t-s}Q_s$ is $1$). It is not a new family.

\begin{theorem}[closed-form table]\label{thm:table}
Write $P/D=e_S\,\sigma$ with $\sigma(z)=\prod_{m\in S}\dfrac{z-1/w_m}{z-w_m}$, and let
$\sigma=\sum_{j\ge0}d_jz^{j}$ at $0$ and $\sigma=\sum_{k\ge0}c_kz^{-k}$ at $\infty$, so
$d_0=\sigma(0)=e_S^{-2}$, $c_0=1$ and $c_1=p_1(w)-p_1(w^{-1})$. Then
\[
  I_a=\begin{cases}
    e_S\,c_{a+1}, & a\ge0,\\[2pt]
    e_S-e_S^{-1}, & a=-1,\\[2pt]
    -e_S\,d_{-a-1}, & a\le-2.
  \end{cases}
\]
\end{theorem}

\begin{proof}
$P(z)=\prod w_m(z-1/w_m)=e_S\prod(z-1/w_m)$, so $P/D=e_S\sigma$ and
$\sigma(0)=\prod(-1/w_m)/(-w_m)=e_S^{-2}$, $\sigma(\infty)=1$. By Lemma~\ref{lem:Phi},
$I_a=e_S\bigl(-\Res_0[z^a\sigma]-\Res_\infty[z^a\sigma]\bigr)$. Now
$\Res_0[z^a\sigma]=d_{-a-1}$, which vanishes unless $a\le-1$; and
$\Res_\infty[z^a\sigma]=-[z^{-1}](z^a\sigma)=-c_{a+1}$, which vanishes unless $a\ge-1$.
The three cases follow, the middle one as
$e_S(-d_0+c_0)=e_S(1-e_S^{-2})=e_S-e_S^{-1}$. The value of $c_1$ comes from
$\sigma=\prod(1-\tfrac{1}{w_mz})\prod(1-\tfrac{w_m}{z})^{-1}
=1+\bigl(\sum w_m-\sum w_m^{-1}\bigr)z^{-1}+O(z^{-2})$.
\end{proof}

\begin{corollary}[the $\lambda=(2,3)$ tie is the member $a=-1$]
$\Phi[P/z]=e_S-1/e_S$.
\end{corollary}

\begin{proposition}[degeneracy is a locus, not a member]\label{prop:deg}
No member $I_a$ vanishes identically in $w$. But if the multiset $\{w_m\}_{m\in S}$ is
self-reciprocal then $e_S=\pm1$ is forced, $c_1=0$, and in particular $I_{-1}\equiv0$ and
$I_0\equiv0$ on that locus.
\end{proposition}

\begin{proof}
$I_0=e_Sc_1=e_S(p_1(w)-p_1(w^{-1}))$ is not the zero function; similarly
$I_{-1}=e_S-e_S^{-1}\not\equiv0$, and for general $a$ the relevant coefficient of
$\sigma$ is a nonzero Laurent polynomial. If $\{w_m\}=\{1/w_m\}$ as multisets then
$e_S=\prod w_m=\prod w_m^{-1}=e_S^{-1}$, so $e_S^2=1$; and $p_1(w)=p_1(w^{-1})$ gives
$c_1=0$. Hence both $I_{-1}=e_S-e_S^{-1}$ and $I_0=e_Sc_1$ vanish there.
\end{proof}

\begin{remark}
Proposition~\ref{prop:deg} is why a sample point must be asserted non-reciprocal: on the
reciprocal locus the \emph{positive control} $I_{-1}=e_S-1/e_S$ reads $0=0$, i.e.\ a
degenerate sample would present as a failing control. The harness asserts
$w_k$ distinct, $w_k\neq0$, $|w_k|\neq1$, $w_jw_m\neq1$, $\{w\}\neq\{1/w\}$ and
$e_S\neq\pm1$.
\end{remark}

\begin{theorem}[reflection duality]\label{thm:dual}
$\sigma_{1/w}=1/\sigma_w$ and $\sigma_w(1/z)=\dfrac{1}{e_S^2\,\sigma_w(z)}$. Consequently
\[
  I_a(1/w)\;=\;-\,I_{-2-a}(w)\qquad\text{for all }a\in\mathbb{Z},
\]
so $(a,w)\mapsto(-2-a,1/w)$ is an involution of the identity family whose unique fixed
index is $a=-1$, on which it acts by $-1$.
\end{theorem}

\begin{proof}
$\sigma_{1/w}(z)=\prod\frac{z-w_m}{z-1/w_m}=1/\sigma_w(z)$ is immediate. Next
\[
  \sigma_w(1/z)=\prod_{m\in S}\frac{\tfrac1z-\tfrac1{w_m}}{\tfrac1z-w_m}
  =\prod_{m\in S}\frac{z-w_m}{w_m(w_mz-1)}
  =\frac{D(z)}{e_S\,P(z)}=\frac{D(z)}{e_S^2\prod(z-1/w_m)}
  =\frac{1}{e_S^2\,\sigma_w(z)} .
\]
Since $e_{S}(1/w)=1/e_S$, Theorem~\ref{thm:table}'s first line of proof gives
\[
  I_a(1/w)=e_S^{-1}\bigl(-\Res_0[z^a\sigma_{1/w}]-\Res_\infty[z^a\sigma_{1/w}]\bigr)
          =e_S^{-1}\bigl(-\Res_0[z^a/\sigma_w]-\Res_\infty[z^a/\sigma_w]\bigr),
\]
and $1/\sigma_w(z)=e_S^2\,\sigma_w(1/z)$, so
$I_a(1/w)=e_S\bigl(-\Res_0[G]-\Res_\infty[G]\bigr)$ with $G(z)=z^a\sigma_w(1/z)$.
Put $F(z)=z^{-a-2}\sigma_w(z)$. Then $G(1/u)/u^2=u^{-a-2}\sigma_w(u)=F(u)$, so
$\Res_\infty[G]=-\Res_{u=0}[G(1/u)/u^2]=-\Res_0[F]$; and
$F(1/u)/u^2=u^{a}\sigma_w(1/u)=G(u)$, so $\Res_\infty[F]=-\Res_0[G]$. Therefore
\[
  -\Res_0[G]-\Res_\infty[G]=\Res_\infty[F]+\Res_0[F]
  =-\bigl(-\Res_0[F]-\Res_\infty[F]\bigr),
\]
and multiplying by $e_S$ gives $I_a(1/w)=-I_{-2-a}(w)$. The index map $a\mapsto-2-a$ is an
involution of $\mathbb{Z}$ with unique fixed point $a=-1$; there it reads
$I_{-1}(1/w)=-I_{-1}(w)$, consistent with $e_S\mapsto e_S^{-1}$ in
$I_{-1}=e_S-e_S^{-1}$.
\end{proof}

\begin{remark}
Theorem~\ref{thm:dual} relates the family at $w$ to the family at $1/w$. It does
\emph{not} give a linear relation among $\{I_a(w)\}_{a}$ at a fixed $w$, and must not be
read as halving the family: it halves the family only after quotienting the parameter
space by $w\mapsto1/w$. The fixed index $a=-1$ is the seam, and it is exactly the
$\lambda=(2,3)$ tie --- which is the structural reason that tie is the distinguished one.
\end{remark}

\section{Head collapse: every head reduces to length one}

The next result needs no residues; it is the expansion of $\rho$ as a polynomial in the
head variable. It applies to \emph{all} head numbers $r$, so it also answers the lift
question Q551 for this component.

\begin{theorem}[head collapse]\label{thm:collapse}
For all $\kappa\in\mathbb{Z}_{\ge1}^r$, $\theta\in\mathbb{Z}^r$ and finite $M$,
\[
  \PsiC_{\kappa,\theta}(M)
  =\sum_{i_1=0}^{\kappa_1-1}\cdots\sum_{i_r=0}^{\kappa_r-1}
    \Biggl(\prod_{j=1}^{r}(-1)^{\kappa_j-1-i_j}\,
      e_{\kappa_j-1-i_j}\bigl(w_1,\dots,w_{\kappa_j-1}\bigr)\Biggr)
    \PsiC_{(1^r),\,\theta'}(M),
  \qquad \theta'_j=\theta_j-\kappa_j+1+i_j .
\]
In particular, for one head,
$\PsiC_{(s),t}(S)=\sum_{i=0}^{s-1}(-1)^{s-1-i}e_{s-1-i}(w_1,\dots,w_{s-1})\,
\PsiC_{(1),\,t-s+1+i}(S)$.
\end{theorem}

\begin{proof}
In \eqref{eq:Psidef} the only dependence on $\kappa$ and $\theta$ is through the factor
$\prod_j\rho(\kappa_j,X_j,w)X_j^{\theta_j}$. For each $j$,
\[
  \rho(\kappa_j,X_j,w)=\prod_{k=1}^{\kappa_j-1}\Bigl(1-\frac{w_k}{X_j}\Bigr)
  =X_j^{1-\kappa_j}\prod_{k=1}^{\kappa_j-1}(X_j-w_k)
  =X_j^{1-\kappa_j}Q_{\kappa_j}(X_j),
\]
so $\rho(\kappa_j,X_j,w)X_j^{\theta_j}=X_j^{\theta_j-\kappa_j+1}Q_{\kappa_j}(X_j)$.
Expanding $Q_{\kappa_j}(X)=\prod_{k<\kappa_j}(X-w_k)
=\sum_{i=0}^{\kappa_j-1}(-1)^{\kappa_j-1-i}e_{\kappa_j-1-i}(w_1,\dots,w_{\kappa_j-1})X^{i}$
gives
\[
  \rho(\kappa_j,X_j,w)X_j^{\theta_j}
  =\sum_{i_j=0}^{\kappa_j-1}(-1)^{\kappa_j-1-i_j}
     e_{\kappa_j-1-i_j}(w_1,\dots,w_{\kappa_j-1})\,X_j^{\theta_j-\kappa_j+1+i_j}.
\]
Multiplying over $j$, substituting into \eqref{eq:Psidef}, and exchanging the finite sums
with the head sum (every other factor in \eqref{eq:Psidef} is independent of
$\kappa,\theta$) yields the stated expansion, each term being the definition of
$\PsiC_{(1^r),\theta'}$ with $\theta'_j=\theta_j-\kappa_j+1+i_j$.
\end{proof}

\begin{corollary}
Theorem~\ref{thm:collapse} with $r=1$, combined with \eqref{eq:Ia} and
Theorem~\ref{thm:table}, is a closed form for $\PsiC_{(s),t}(S)$ for \emph{every}
$(s,t)$:
\[
  e_S\,\PsiC_{(s),t}(S)
  =\sum_{i=0}^{s-1}(-1)^{s-1-i}e_{s-1-i}(w_1,\dots,w_{s-1})\,I_{\,t-s+i},
\]
with each $I_b$ given by the three-case table. This is strictly stronger than
\eqref{eq:proved}, which still required two residues of an $s$-dependent function. The
head data $(w_1,\dots,w_{s-1})$ enters only through fixed elementary symmetric
polynomials and is entirely decoupled from $S$.
\end{corollary}

\section{The head sum is an alternant, and the $\rho$-shifted vanishing law}

\begin{theorem}[alternant expansion]\label{thm:alt}
Let $m_j=\theta_j+j-1$ for $j=1,\dots,r$. For every finite $M$ and $\theta\in\mathbb{Z}^r$,
\[
  \PsiC_{(1^r),\theta}(M)
  =\sum_{\substack{T\subseteq M\\ |T|=r}} B(T)\cdot
    \frac{\det\bigl(x_i^{\,m_j}\bigr)_{i,j=1}^{r}}{\prod_{i<k}(x_k-x_i)},
\]
where $x_1,\dots,x_r$ are the $w$-values of $T$ in increasing index order and
\[
  B(T)=\prod_{x\in T}\bigl(1-x^{-2}\bigr)\cdot
       \prod_{\{x,y\}\subseteq T}\Bigl(1-\tfrac{1}{xy}\Bigr)\cdot
       \prod_{t\in M\setminus T}\prod_{x\in T}\frac{1-\frac{1}{w_tx}}{1-\frac{w_t}{x}} .
\]
\end{theorem}

\begin{proof}
Group the sum \eqref{eq:Psidef} (with every $\kappa_j=1$, so $\rho\equiv1$) by the
underlying $r$-subset $T=\{h_1,\dots,h_r\}$. All factors of \eqref{eq:Psidef} except
$\prod_j X_j^{\theta_j}$ and $\prod_{j<k}(1-X_j/X_k)^{-1}$ are symmetric functions of
$T$, and their product is $B(T)$. For a fixed $T$ with values $x_1,\dots,x_r$, the
orderings of $T$ are indexed by $\pi\in\mathfrak{S}_r$ via $X_j=x_{\pi(j)}$, and
\[
  \prod_{j<k}\Bigl(1-\frac{x_{\pi(j)}}{x_{\pi(k)}}\Bigr)
  =\frac{\prod_{j<k}\bigl(x_{\pi(k)}-x_{\pi(j)}\bigr)}{\prod_{k}x_{\pi(k)}^{\,k-1}}
  =\operatorname{sgn}(\pi)\,\frac{\prod_{i<k}(x_k-x_i)}{\prod_k x_{\pi(k)}^{\,k-1}},
\]
the last step being the Vandermonde sign rule. Hence the ordered contribution of $T$ is
\[
  \sum_{\pi\in\mathfrak{S}_r}
    \frac{\prod_j x_{\pi(j)}^{\theta_j}\,\prod_k x_{\pi(k)}^{k-1}}
         {\operatorname{sgn}(\pi)\prod_{i<k}(x_k-x_i)}
  =\frac{1}{\prod_{i<k}(x_k-x_i)}
   \sum_{\pi\in\mathfrak{S}_r}\operatorname{sgn}(\pi)\prod_{j=1}^{r}x_{\pi(j)}^{\,\theta_j+j-1},
\]
using $\operatorname{sgn}(\pi)^{-1}=\operatorname{sgn}(\pi)$. The last sum is
$\det(x_i^{m_j})_{i,j}$ with $m_j=\theta_j+j-1$, by the permutation expansion of the
determinant. Multiplying by $B(T)$ and summing over $T$ gives the claim.
\end{proof}

\begin{corollary}
When the $m_j$ are pairwise distinct, sorting them as $m_{\mathrm{sort}}=\mu+(0,1,\dots,r-1)$
makes $\det(x_i^{m_j})/\prod_{i<k}(x_k-x_i)=\pm s_\mu(T)$ a (Laurent-)Schur polynomial, so
$\PsiC_{(1^r),\theta}$ is a $B$-weighted sum of Schur polynomials over the $r$-subsets of
$M$. Combined with Theorem~\ref{thm:collapse}, \emph{every} $\PsiC_{\kappa,\theta}$ is such
a sum.
\end{corollary}

\begin{theorem}[$\rho$-shifted vanishing law]\label{thm:vanish}
If $\theta_j-\theta_k=k-j$ for some $j<k$, then $\PsiC_{(1^r),\theta}(M)=0$ identically in
$w$, for every $M$.
\end{theorem}

\begin{proof}
$\theta_j-\theta_k=k-j$ is exactly $m_j=m_k$ for $m_j=\theta_j+j-1$. Then the matrix
$(x_i^{m_j})$ has two equal columns, so its determinant is zero for every $T$, and every
summand of Theorem~\ref{thm:alt} vanishes.
\end{proof}

\begin{remark}[direction of the vanishing law]\label{rem:converse}
The converse --- that $\PsiC_{(1^r),\theta}\neq0$ when the $m_j$ are pairwise distinct ---
is \textbf{not} proved here. Each summand is then a nonzero $\pm B(T)s_\mu(T)$, but the sum
over $T$ could in principle cancel. It is \emph{measured} over the full boxes
$\theta\in[-4,4]^2$ ($81$ points) and $\theta\in[-3,3]^3$ ($343$ points), where the
predicted zero set and the measured zero set agree \emph{as sets, in both directions},
with no $M$-dependence. I grade the forward direction \texttt{proved} and the converse
\texttt{computed}.
\end{remark}

\section{Verification}\label{sec:verif}

All computations are exact: \texttt{fractions.Fraction} throughout, with SymPy used only
for symbolic residues and only via \texttt{cancel}/\texttt{simplify}
(\texttt{nsimplify} is never called --- it has twice manufactured false mismatches by
returning irrational forms for rationals). SageMath is not installed in this container;
the two routes below --- a SymPy residue computation and a pure-\texttt{Fraction}
head-sum --- reach the same numbers by different mechanisms, so they are independent
corroboration rather than one engine read twice. Scripts:
\texttt{scratch/2026-10-22-c2-Q553/\{q553,q553b,q553c,q553d,vanish,vanish2,alternant\}.py}.

\begin{center}
\begin{tabular}{llr}
\hline
Claim & Range & Result\\
\hline
Thm~\ref{thm:carrier}, set equality both directions & $|S|=3,4$; $s\le3$; $-2\le t\le2$ & $60/60$\\
Thm~\ref{thm:one}, $\sum_{\mathbb{P}^1}\Res=0$ & $|S|=3,4$; $s\le2$; $t\in\{-2,0,2\}$ & $18/18$\\
Thm~\ref{thm:table} vs.\ Lagrange definition & $5$ sets $S$, $-4\le a\le4$ & $45/45$\\
$I_{-1}=e_S-1/e_S$ (positive control) & $5$ sets, with and without $1\in S$ & $5/5$, fires\\
Prop.~\ref{prop:deg}, vanishing members, generic $w$ & as above & $0$ of $45$\\
Prop.~\ref{prop:deg}, reciprocal locus $\Rightarrow I_0=0$ & $n=2,4$ & confirmed\\
Thm~\ref{thm:dual}, $\sigma_w(1/z)=1/(e_S^2\sigma_w)$ & $4$ sets & $4/4$\\
Thm~\ref{thm:dual}, $I_a(1/w)=-I_{-2-a}(w)$ & $4$ sets, $-4\le a\le3$ & $32/32$\\
Cor.\ to Thm~\ref{thm:collapse}, closed form all $(s,t)$ & $8$ sets; $s\le4$; $5$ values of $t$ & $160/160$\\
Thm~\ref{thm:collapse}, $r=1$, pure \texttt{Fraction} & $10$ sets; $s\le5$; $-4\le t\le3$ & $400/400$\\
Thm~\ref{thm:collapse}, $r=2$ & $5$ sets, $7$ $\kappa$, $5$ $\theta$ & $175/175$\\
Thm~\ref{thm:collapse}, $r=3$ & $3$ sets, $4$ $\kappa$, $3$ $\theta$ & $36/36$\\
Thm~\ref{thm:alt}, $r=1,2,3$ & $15+28+18$ cases & $61/61$\\
Thm~\ref{thm:vanish}, zero set as a set, $r=2$ & $\theta\in[-4,4]^2$, $3$ sets $M$ & $81/81$\\
Thm~\ref{thm:vanish}, zero set as a set, $r=3$ & $\theta\in[-3,3]^3$, $2$ sets $M$ & $343/343$\\
\hline
\multicolumn{3}{l}{\emph{Negative controls (each must read $0/N$):}}\\
drop the sign $(-1)^{\kappa_j-1-i_j}$ in Thm~\ref{thm:collapse} & live samples only & $0/20$\\
shift the head window to $w_2,\dots,w_s$ & $r=1$, $s\le4$ & $0/18$\\
drop the $+j$ shift in $m_j$ (Thm~\ref{thm:alt}) & $r=2$ & $0/8$\\
\hline
\end{tabular}
\end{center}

\begin{remark}[a half-dead control was the discovery instrument]
The first run of the sign-dropping control read $4/12$ matching, not $0/12$. Diagnosis
(rather than hand-waving) showed all four ``matches'' were \emph{degenerate samples}: the
sign-flipped term was $\PsiC_{(1,1),(0,-1)}$, which is identically zero. That accident is
what exposed Theorem~\ref{thm:vanish}, and $(0,-1)$ is the $r=2$ instance of
$\theta_j-\theta_k=k-j$. Re-run on samples chosen off the vanishing wall, the control reads
$0/20$. My first guess at the locus --- $\theta_1+\theta_2=-1$ --- was \emph{wrong}, and
the datum that killed it was $\theta=(0,0,-2)$ at $r=3$, which vanishes with no pair
difference equal to $1$; the $\rho$-shift $m_j=\theta_j+j-1$ is what fits, $424/424$.
\end{remark}

\section{T0: the Schur-Q kernel identification (retraction and result)}\label{sec:T0}

\textbf{Retraction, 2026-10-23.} The text of this section as first written graded a
collaborator's note at commit \texttt{c53c581} as resting on an \emph{unreadable source}, on
the evidence that \texttt{git cat-file -t c53c581} returned \texttt{fatal: Not a valid object
name} in all eight clones of his repository on this machine. \textbf{That grade was wrong and
the hash is valid.} Measured on 2026-10-23 in both directions: \texttt{cat-file} returns
\texttt{fatal} \emph{before} \texttt{git fetch} and \texttt{commit} \emph{after}; \texttt{gh
api} resolves \texttt{c53c5813c1914ceca8d682b101d2e3e4f0b04777}, dated 2026-10-22T12:41:15Z.
All eight clones shared a single stale fetch, so they were eight readings of one instrument and
not eight instruments; counting clones is not replication, and the authority for a remote
object is the remote. The collaborator's future-dated filename is his standing naming
convention, not a tell, and that inference is withdrawn too.

The verdict was logged as \textbf{NO MATCH on unreadable source}, explicitly \emph{not} on a
measured mismatch. That is what makes this a retraction and not a withdrawal: no mathematics in
this file depended on it.

\textbf{And the identification holds.} With the source now read, the collaborator's iterated
residue formula couples its variables through $K(\tau,t)=(\tau-t)/(\tau t-1)$, and under his
Cayley map $t=(1+z)/(1-z)$ this becomes the Schur-Q kernel $(z_r-z_s)/(z_r+z_s)$ (verified
symbolically). Writing $\sigma_w(z)=\prod_{m\in S}(z-1/w_m)/(z-w_m)=P(z)/(e_SD(z))$ as in
\S\ref{sec:basis}, the comparison is not an analogy but an identity:
\[
  \prod_{m\in S}K(z,w_m)\;=\;e_S\,\sigma_w(1/z)\;=\;\frac{1}{e_S\,\sigma_w(z)} .
\]
Verified symbolically for $|S|\le5$, with five controls all firing (dropping $e_S$, using
$e_S^2$, the two wrong kernels $(a-b)/(ab+1)$ and $(a+b)/(ab-1)$, and omitting the reflection).
So his coupling kernel, multiplied over the tail set, \emph{is} the reflection operator of
Theorem~\ref{thm:dual}, and Theorem~\ref{thm:dual}'s functional equation
$\sigma_w(1/z)=1/(e_S^2\sigma_w(z))$ is exactly the functional equation of his kernel product.

One correction in the other direction: the kernel is \emph{not} antisymmetric under
$t\mapsto1/t$ in a single variable, as I had expected. It is reciprocal-valued there
($K\mapsto1/K$) and invariant under the simultaneous inversion; its antisymmetry is under the
\emph{swap} $r\leftrightarrow s$, a different group element. Any argument that pairs the two
involutions by their fixed points must use the simultaneous form.

\section{Gaps and open points}\label{sec:gaps}

\begin{enumerate}
\item \textbf{Converse of Theorem~\ref{thm:vanish}} --- non-vanishing off the walls is
      \texttt{computed} ($424/424$), not proved; see Remark~\ref{rem:converse}. A proof
      presumably needs a positivity or leading-term argument on
      $\sum_T B(T)s_\mu(T)$, and $B(T)$ is not of one sign.
\item \textbf{Generic clauses in Theorem~\ref{thm:carrier}} --- the memberships of $0$
      for $t<s-1$ and of $\infty$ for $t\ge1$ are coefficient non-vanishing statements,
      true off a proper subvariety; the set equality is verified at $60$ random rational
      points but is a generic, not universal, claim. The unconditional cases are
      $t=s-1$, $t=0$, and all $w_m$ with $m\ge s$.
\item \textbf{Q553's counting question is answered in the negative but reopened in a
      better form.} There is no finite ``number of admissible groupings''; the identity
      family is $\mathbb{Z}$-indexed with an explicit basis and an explicit closed form.
      Whether $\{I_a\}_{a\in\mathbb{Z}}$ is linearly independent \emph{as functions of
      $w$} --- which is the only sense in which a count is meaningful --- is not settled
      here.
\item \textbf{T0} --- a quoted peer artifact must be stamped with a commit hash that has
      been resolved by \texttt{git cat-file} \emph{in the same action as the quote}. Here
      the hash was present and false. A locator written inside my own brief is not a
      locator: the quoting artifact and the quoted artifact were produced by the same slot
      minutes apart, so the quote's only corroboration was the brief that needed it.
\end{enumerate}

\end{document}
